\documentclass[10pt]{article}
\usepackage[margin=20mm]{geometry}
\usepackage{amsmath,amssymb,mathtools,fontspec,unicode-math,xurl,hyperref}
\IfFontExistsTF{DejaVu Serif}{\setmainfont{DejaVu Serif}}{\setmainfont{Latin Modern Roman}}
\setmathfont{latinmodern-math.otf}
\setlength{\parindent}{0pt}\setlength{\parskip}{6pt}\sloppy
\title{Initial scientific review of
Sharp bounds on the Nusselt number in Rayleigh–Bénard convection and a bilinear estimate by Coifman–Meyer}
\author{}
\date{26 October 2026}
\begin{document}
\maketitle
Abstract: Proposition 2.3 is false with the stated arbitrary dependence of the second input on the scale variable. A family of real, uniformly bounded, spatially compact inputs satisfies every hypothesis, while the bilinear output has an \(L^2\) norm growing at least as the square root of the number of selected scales and both right-hand norms remain bounded. The appendix’s treatment of products with zero output frequency does not justify the required square-function bound. The proposition is used directly in the core heat-flux proof, so this is one Major finding. No counterexample to a solution of the coupled convection equations or to the final physical upper bound is claimed. A replacement argument has not been verified.\par
\section*{Identity, version and coverage}
The assigned article is Sagun Chanillo and Andrea Malchiodi, Sharp bounds on the Nusselt number in Rayleigh–Bénard convection and a bilinear estimate by Coifman–Meyer, M25-0066, Inventiones mathematicae 240 (2025), 633–660, DOI 10.1007/s00222-025-01326-z [paper]. The reviewed original is the acquired publisher version of record, 28 pages. Printed page \(p\) corresponds to one-based PDF page \(p-632\).\par
All 1561 lines of the native PDF extraction were read continuously, including the full appendix and references. Twenty-one original PDF pages were visually inspected: \[1,2,3,4,5,6,7,8,9,12,13,15,16,17,18,20,21,23,24,25,26.\] This includes every figure page, 8, 13, 18 and 24. The exact hypotheses of Proposition 2.3, its application on p. 648, and the appendix on pp. 656–658 were checked in the page images. Selected passages of the supplied arXiv:2001.01662v4 TeX were inspected statically as related author-source evidence, without treating that earlier manuscript as the journal source. All nine manifest-listed files were hash checked.\par
The principal audit followed the normalization, boundary conditions, Green-kernel reduction, localization, bilinear estimate and final dependency chain. It did not independently reconstruct all cited PDE regularity results, all asymptotic derivative calculations, or every Fourier correction formula. No author code was executed, and no fluid simulation was performed.\par
\section*{Scope of the claimed theorem and the analytic input}
The physical setting is the infinite-Prandtl-number Boussinesq system with horizontal periodicity, no-slip velocity at the two plates, and fixed temperatures one and zero. The rescaled height is \(H=\mathrm{Ra}^{1/3}\), with \(H\gg1\). Theorem 1.2, p. 635, claims a uniform bound on the rescaled Nusselt number; Corollary 1.3 converts it to \(\mathrm{Nu}_{\mathrm{phys}}\le C_0\mathrm{Ra}^{1/3}\). The theorem does not cover arbitrary finite Prandtl number or supply the numerical constant discussed in the introduction.\par
The temperature fluctuation \(\theta\) has zero horizontal mean and \(|\theta|\le1\) under the maximum-principle bound. The vertical velocity is obtained from a clamped fourth-order equation. The proof separates the forcing into lower and upper half-strips, represents the lower part by a half-space kernel, and corrects its upper boundary values by a biharmonic function. The near-bottom correlation bound is the essential point: Proposition 3.2 invokes Proposition 2.3 to eliminate dependence on the number of scales between the boundary and the height \(H\).\par
The finding below concerns the independently stated functional inequality. Its hypotheses require no heat equation, no harmonicity of the second input in the scale variable, and no common underlying function across scales. Those distinctions are decisive.\par
\section*{Problem 1: the scale-dependent bilinear estimate is false}
\subsection*{Source claim and location}
Equation reference: claimed. Proposition 2.3, printed p. 640 (PDF p. 8), uses \[P_t(x)=\frac{1}{2\pi}\frac{t}{(t^2+|x|^2)^{3/2}},
 \qquad Q_tf=\psi_t*f=t\nabla(P_t*f).\] For \(\rho\in\mathrm{BMO}(\mathbb R^2)\), bounded compactly supported \(F\) on \(\mathbb R^2\times(0,\infty)\), and measurable \(|m|\le1\), it asserts \[
 \left\|\int_0^\infty Q_t\rho\cdot Q_tF^t\,m(t)\,\frac{dt}{t}\right\|_2
 \le C\|\rho\|_{\mathrm{BMO}}\,\|N(Q_tF^t)\|_2 ,\] with \(C\) independent of all inputs. Here \(NF(x)=\sup_{|z-x|<t}|F(z,t)|\), as defined in (17), p. 639. The paragraph below Remark 2.4 explicitly identifies arbitrary \(t\)-dependence as the extension beyond the original setup.\par
\subsection*{Construction and verification of every hypothesis}
Equation reference: inputs. Use angular frequencies \(q_j=5^j\), \(1\le j\le N\), and disjoint intervals \[I_j=[q_j^{-1},2q_j^{-1}].\] Let \[E=[0,2\pi]^2,\qquad Q=[-2\pi,4\pi]^2,\qquad
 g_N(x_1)=\sum_{j=1}^N\cos(q_jx_1),\qquad
 r_N(x_1)=\operatorname{sgn}g_N(x_1).\] The value at zeros may be set to zero. Define \[
 \rho_N(x)=\mathbf1_Q(x)r_N(x_1),\qquad
 F_N(x,t)=
 \begin{cases}
 \mathbf1_Q(x)\cos(q_jx_1),&t\in I_j,\\
 0,&t\notin\bigcup_{j=1}^NI_j,
 \end{cases}
 \qquad m(t)=1.\] These are real measurable functions. Both inputs have absolute value at most one, and \(\|\rho_N\|_{\mathrm{BMO}}\le2\). In fact \(\rho_N\) is spatially compact as well, although the proposition does not require that. The support of \(F_N\) lies in \[Q\times[5^{-N},2/5],\] a compact subset of the open upper half-space for each finite \(N\). Its approach to \(t=0\) as \(N\) grows violates no stated hypothesis: the claimed constant is independent of \(F\).\par
Equation reference: kernel, maximal. The vector kernel is \[
 \psi_t(y)=-\frac{3}{2\pi}\frac{t^2y}{(t^2+|y|^2)^{5/2}},
 \qquad \|\psi_t\|_{L^1(\mathbb R^2;\mathbb R^2)}=1.\] Consequently \(|Q_tF_N^t|\le1\). It vanishes for \(t\) outside the selected intervals. Write \(a=(\pi,\pi)\). When \(|x-a|\) is larger than a fixed radius containing \(Q\), \(t\le2/5\) and \(|z-x|<t\) imply \[|Q_tF_N^t(z)|\le C_Q|x-a|^{-4}.\] Indeed, the integration is over the fixed square \(Q\) and [kernel] (\#kernel) is bounded by a constant times \(t^2|z-y|^{-4}\) there. Combining this with the bound one on a fixed disk gives \[
 \|N(Q_tF_N^t)\|_{L^2(\mathbb R^2)}\le C_Q,\] uniformly in \(N\). Thus the entire right-hand side of [claimed] (\#claimed) is bounded by \(2CC_Q\).\par
\subsection*{Exact lower bound for the left-hand side}
Equation reference: mean. Let \(\mathbb E\) denote normalized integration over \(E\) for functions periodic in \(x_1\) and constant in \(x_2\). Temporarily omit both spatial cutoffs and write \(\overline B_N\) for the corresponding bilinear expression with first input \(r_N\) and second input \(\cos(q_jx_1)\) on \(I_j\). For the paper’s Poisson normalization, \[P_t*\cos(q_jx_1)=e^{-tq_j}\cos(q_jx_1).\] Periodic integration, or the skew-adjointness of horizontal differentiation, therefore gives \[\mathbb E\big[Q_tr_N\cdot Q_t\cos(q_jx_1)\big]
 =(tq_j)^2e^{-2tq_j}\mathbb E[r_N\cos(q_jx_1)].\] Hence \[
 \mathbb E\overline B_N
 =A\,\mathbb E[r_Ng_N]
 =A\,\mathbb E|g_N|,
 \qquad
 A=\int_1^2u e^{-2u}\,du=\frac{3e^{-2}-5e^{-4}}4>0.\] The use of the full vector kernel causes no cancellation here: the periodic comparison functions depend only on \(x_1\), so their second derivative components vanish.\par
Equation reference: L1. There are no non-pairing signed four-term relations among the frequencies \(5^j\). To see this, a largest frequency occurring once cannot be balanced by three smaller ones; two occurrences with the same sign cannot be balanced by two smaller ones; and three occurrences cannot be balanced by one smaller frequency. The remaining possibilities are opposite pairs. Orthogonality thus yields the exact moments \[\mathbb Eg_N^2=\frac N2,\qquad
 \mathbb Eg_N^4=\frac{3N^2}{4}-\frac{3N}{8}\le\frac{3N^2}{4}.\] Hölder interpolation between \(L^1,L^2,L^4\) gives \[
 \mathbb E|g_N|\ge
 \frac{(\mathbb Eg_N^2)^{3/2}}{(\mathbb Eg_N^4)^{1/2}}
 \ge\sqrt{\frac N6}.\]\par
Equation reference: divergence. It remains to verify that compact support does not invalidate this lower bound. For \(x\in E\), the distance to \(Q^c\) is at least \(d=2\pi\). From [kernel] (\#kernel), \[\int_{|y|\ge d}|\psi_t(y)|\,dy\le \frac{3t^2}{2d^2}
 =\frac{3t^2}{8\pi^2}.\] Each cut input differs after convolution from its uncut counterpart by at most this amount. All four convolved factors have norm at most one. The difference of their scalar products is consequently at most twice this tail. If \(B_N\) denotes the actual expression with the inputs [inputs] (\#inputs), then \[|B_N(x)-\overline B_N(x)|
 \le\frac{3}{4\pi^2}\sum_{j=1}^N\int_{I_j}t^2\frac{dt}{t}
 \le\frac{3}{64\pi^2}\quad(x\in E),\] because \(\sum_{j\ge1}5^{-2j}=1/24\). By [mean] (\#mean)–[L1] (\#L1) and Cauchy–Schwarz on \(E\), \[
 \|B_N\|_{L^2(\mathbb R^2)}
 \ge |E|^{1/2}\left(A\sqrt{\frac N6}-\frac{3}{64\pi^2}\right).\] This tends to infinity, whereas [maximal] (\#maximal) and the BMO bound keep the claimed right-hand side bounded. No universal \(C\) in Proposition 2.3 exists.\par
This is an analytic counterexample family, not an extrapolation from numerical samples. Independent exact Fourier-moment checks for thirteen values of \(N\), up to \(64\), corroborate the moment calculation. The output records \(A=0.0786069138165418\). Small-grid quadrature is explicitly labeled illustrative.\par
\subsection*{Where the appendix’s extension fails}
The appendix distinguishes the separated input \(F(x,t)=f(x)m(t)\) in (55), p. 657, from the more general input it wants to admit. The support argument preceding (59) is valid for a product whose output spectrum stays away from zero. That does not cover the interaction of opposite frequencies in the present family.\par
The multiplier \(\widehat\psi_3\) on p. 657 is allowed to be supported in a ball or a proper annulus. On p. 658, immediately after (59), the first factor is bounded by a constant times \(\|h\|_2\): \[\left(\int_0^\infty\|\psi_{3,t}*h\|_2^2\,\frac{dt}{t}\right)^{1/2}
 \lesssim\|h\|_2.\] If the multiplier is one near zero, then \(\psi_{3,t}*h\to h\) in \(L^2\) as \(t\downarrow0\), and this integral diverges for every nonzero \(h\in L^2\). An annular multiplier avoids that divergence but cannot preserve a product’s nonzero constant-frequency contribution. The general reduction must control this resonant case by additional information; merely retaining the support inclusion does not do so. The counterexample proves that the proposed final extension cannot hold.\par
\subsection*{Counterevidence, impact and repair boundary}
Several possible restrictions do not remove the example. Both inputs are bounded by one, \(\rho_N\) is also compactly supported, all scale integrals are over a finite interval bounded away from zero for each \(N\), and the non-tangential maximal function has a uniform integrable tail. The proposition assumes no smoothness. If desired, for each fixed \(N\) one can smooth the spatial and temporal cutoffs and the bounded sign function, retaining uniform bounds and passing to the displayed finite-scale integrals; the obstruction is not a singularity of the inputs.\par
The kernel used in the application on p. 648 is \(t\nabla P_{t/2}\). The same argument works for it: the positive scale coefficient becomes \(\int_1^2u e^{-u}\,du=2e^{-1}-3e^{-2}\), and all cutoff and maximal-function constants remain uniform. Changing this normalization cannot repair the estimate.\par
The core dependency is explicit: \[\begin{aligned}
 &\text{Proposition 2.3}\ \longrightarrow\
 \text{(37), Proposition 3.2}\\
 &\hspace{2em}\longrightarrow\
 \text{Proposition 3.3}\ \longrightarrow\
 \text{Theorem 1.2 and Corollary 1.3}.
\end{aligned}\] Equation (37), p. 648, is the direct invocation. Proposition 3.3, p. 649, is then used in the proof of Theorem 1.2 on p. 655. The later biharmonic boundary correction does not supply a substitute for the failed bilinear bound. This establishes a Major defect in a core derivation.\par
The family is not asserted to solve the coupled temperature–velocity evolution. It therefore does not disprove the long-time physical Nusselt bound. Restricting the scale dependence to a fixed underlying function would remove the freedom used here, but the temperature slices in the application do not have that form. No corrected proposition sufficient for the application has been proved in this review. A repair needs a different analytic estimate or additional PDE structure that controls the accumulating interactions.\par
\section*{Overall assessment and other checked material}
The Poisson normalization, kernel homogeneity, half-space/strip separation, spatial localization and role of the no-slip boundary correction were followed. The local integrability near the Green-kernel singularity and the distinction between the critical cubic far-field term and its integrable remainder were checked at the level used in Lemma 3.1. These parts do not resolve Problem 1.\par
The report does not elevate isolated sign, coefficient or typesetting discrepancies in the displayed Fourier formulas into additional findings. In particular, zero horizontal mean suppresses the zero mode in the actual forcing, and a global sign convention does not repair or negate an absolute-value bilinear counterexample. An independent rederivation of every boundary-correction coefficient is outside the completed scope.\par
The arXiv record [arxiv] reports that an early version was withdrawn over a scaling argument and a previous Proposition 3.1. The current journal article acknowledges a problem pointed out in a previous version. This historical attribution is retained; it is not evidence that the present counterexample was previously identified or corrected. The acquired v4 source contains the same scale-dependent proposition and the corresponding appendix argument at the inspected passages. The publisher and arXiv records and a limited correction search did not identify a later correction of this claim at consultation; that search was not exhaustive.\par
Item | Disposition\par
Issue judgment | One Major: Proposition 2.3 fails for bounded compact inputs.\par
Supported impact | A required estimate in the core proof is invalid; the physical theorem itself is not refuted here.\par
Repair verification | No replacement sufficient for the application verified.\par
Adoption | Not requested; no author contact.\par
Remaining work | Independent assessment of the counterexample and a valid replacement heat-flux argument.\par
Scientific decision | Pending; human reviews remain empty.\par
9 S. Chanillo and A. Malchiodi, Sharp bounds on the Nusselt number in Rayleigh–Bénard convection and a bilinear estimate by Coifman–Meyer, Invent. Math. 240 (2025), 633–660, DOI: 10.1007/s00222-025-01326-z (https://doi.org/10.1007/s00222-025-01326-z). Assigned publisher PDF. S. Chanillo and A. Malchiodi, Sharp bounds on the Nusselt number in Rayleigh–Bénard convection and a bilinear estimate via Carleson measures, arXiv:2001.01662v4, 11 April 2024. Supplied source and primary version record; related earlier manuscript, not substituted for the journal PDF.\par
\section*{Supporting checks}
The related author TeX SHA-256 is\par
The final independent-check output SHA-256 is\par
The independent Python script `bilinear\_checks.py` verified exact integer/fraction moments for \(N=1,\ldots,10,16,32,64\). Quadrature of the scale integral and sampled \(L^1\) moments are corroboration only. Actual output and source-image bindings are preserved in provenance. Static author-TeX reading covered zero-based lines 626–681, 2078–2138 and 2140–2194; no author code was executed.\par
The full Coifman–Meyer book PDF could not be read: the web service rejected its 14211433-byte size. This failed access is recorded in `literature\_check.json`. The counterexample is independent of that book; its original theorem is not alleged to fail.\par
\section*{Supporting checks}
Its SHA-256 is\par
The receipt binds exact source, compiler, logs, PDF and every inspected page image. Its dependency inventory is a full cache snapshot, a superset of used dependencies; final source/PDF hashes remain external.\par
This assistant performed all science and artifact checks; Work scheduled only. Human reviews are empty; scientific decision and independent acceptance remain pending. Token:\par
\end{document}
